\documentclass[10pt,a4paper]{amsart}
\usepackage[margin=19mm]{geometry}
\usepackage{amsmath,amssymb,mathtools}
\usepackage[scr=rsfs,cal=euler]{mathalfa}
\usepackage{xfrac}
\usepackage[unicode,hidelinks,bookmarks=false]{hyperref}
\newcommand{\ie}{i.e.\ }
\newcommand{\cf}{cf.\ }
\newcommand{\resp}{respectively\ }
\newcommand{\Subsection}{\subsection}
\providecommand{\nobreakdash}{\nobreak\hbox{-}}
\setlength{\emergencystretch}{2em}
\multlinegap0pt
\usepackage{mathtools}
\usepackage{amssymb}
\usepackage[shortlabels]{enumitem}
\newtheorem{theorem}{Theorem}
\newtheorem{assumption}{Assumption}
\newcommand{\E}{\mathrm{E}}
\title{Inference Functionals and Observation Operators \\ for Distributional Statistical Models}
\author{R. Labouriau}
\date{Source: arXiv:2605.19189v1 (CC BY 4.0)}
\begin{document}
\maketitle

\noindent\emph{Source and licence.} Inference Functionals and Observation Operators \\ for Distributional Statistical Models. Version arXiv:2605.19189v1 (CC BY 4.0), distributed by arXiv under the Creative Commons Attribution 4.0 International licence, http://creativecommons.org/licenses/by/4.0/, as recorded in the arXiv licence metadata for this exact version and shown on the abstract page https://arxiv.org/abs/2605.19189v1. Full licence notice: \texttt{licenses/arxiv-2605.19189-CC-BY-4.0.txt}.\par\smallskip

\noindent\textit{Editorial scope.} Counted unit: the article's theorem thm:AN (source0.tex lines 1028--1039). The article's complete original proof of it is delivered with it (source0.tex lines 1041--1054, one \texttt{proof} environment of 14 lines). No mathematics is written, repaired or re-derived: the statement and the proof are the article's own lines, in the article's own notation, and no retained line is substituted or re-flowed.  Retained uncounted as the source-local context the proof needs: lines 946--961; lines 998--1008.  Nothing was omitted from any retained span, so the package carries no omission record.  No retained line contains a citation, so the wrapper carries no bibliography and no bibliography entry.  No retained line contains a figure environment, an image-inclusion command or drawing code, so no image is delivered and the package declares no asset dependency.  Editorial formatting only, in addition: an AMS \texttt{amsart} preamble that restates the article's own declarations and macros used by the retained lines, namely the environments \texttt{theorem}, \texttt{assumption} on separate counters, together with the standard packages those lines need.  The retained environments print in this document as Theorem 1, Assumption 1 in order of appearance, whereas the article prints the same items with the numbering of its own full document.  The complete native source of the article as distributed in the arXiv e-print for this exact version is bundled unchanged as \texttt{sources/200900/source0.tex}.
\begin{theorem}[Consistency]\label{thm:consistency}
Suppose that:
\begin{enumerate}
\item \textbf{Glivenko--Cantelli:} for each compact $K \subset \Theta$,
$\sup_{\theta \in K} \|\bar\Psi_n(\theta) - \bar\Psi(\theta)\|
\xrightarrow{P} 0$.

\item \textbf{Identification:} $\bar\Psi(\theta_0) = 0$ and $\theta_0$ is the unique root of
$\bar\Psi(\theta) = 0$.

\item \textbf{Continuity:} the map $\theta \mapsto \bar\Psi(\theta)$ is
continuous.
\end{enumerate}
Then any sequence of solutions $\hat\theta_n$ of $\bar\Psi_n(\theta) = 0$
satisfies $\hat\theta_n \xrightarrow{P} \theta_0$.
\end{theorem}

\begin{assumption}\label{ass:diff}
For $\theta$ in a neighbourhood of $\theta_0$:
\begin{enumerate}
\item $\theta \mapsto \Psi(y, \theta)$ is differentiable for
$P_{\theta_0}^{\mathcal O}$-almost all~$y$.

\item Differentiation under the expectation is valid:
$\partial_\theta \bar\Psi(\theta)
= \E[\partial_\theta \Psi(Y, \theta)]$.
\end{enumerate}
\end{assumption}

\begin{theorem}[Asymptotic normality]\label{thm:AN}
Under the conditions of Theorem~\ref{thm:consistency},
Assumption~\ref{ass:diff}, $S(\theta_0)$ non-singular, and
$V(\theta_0)$ finite and positive definite,
\[
\sqrt{n}\,(\hat\theta_n - \theta_0)
\;\xrightarrow{\;\mathcal D\;}\;
\mathcal N\!\left(0,\;
S(\theta_0)^{-1}\, V(\theta_0)\, S(\theta_0)^{-\top}
\right).
\]
\end{theorem}

\begin{proof}
Taylor expansion of $\bar\Psi_n$ around $\theta_0$:
\[
0 = \bar\Psi_n(\hat\theta_n)
= \bar\Psi_n(\theta_0)
+ \bigl[\partial_\theta \bar\Psi_n(\tilde\theta_n)\bigr]
  (\hat\theta_n - \theta_0),
\]
for some intermediate $\tilde\theta_n$.  By the CLT,
$\sqrt{n}\,\bar\Psi_n(\theta_0) \Rightarrow \mathcal N(0, V(\theta_0))$,
and by the LLN,
$\partial_\theta \bar\Psi_n(\tilde\theta_n) \xrightarrow{P} -S(\theta_0)$.
Slutsky's theorem gives the result.
\end{proof}

\end{document}
